\documentclass{report}
\usepackage{amsmath,amssymb}
\setlength{\parindent}{0mm}
\setlength{\parskip}{1em}
\begin{document}
\begin{center}
\rule{6in}{1pt} \
{\large Justin, W.L. Wan \\
{\bf Multigrid Method for Solving Elliptic Monge-Ampere Equations Arising from Image Registration}}

University of Waterloo \\ 200 University Avenue West \\ Waterloo \\ ON N2L 3G1 \\ Canada
\\
{\tt jwlwan@uwaterloo.ca}\\
Ashmeen Soin\end{center}

The Monge-Amp\`{e}re equation (MAE) is a nonlinear second order partial
differential equation, which can be written as:
\begin{equation}
u_{xx} u_{yy} - u^2_{xy} = f, \qquad on \ \Omega
\label{eqn:MAE}
\end{equation}
where $f>0$ so that the equation is elliptic with a unique convex solution.
Monge-Amp\`{e}re equations arise in many areas such as differential geometry
and other applications. In image registration, one is interested in
transforming one image to align with another image. One approach is based on
the Monge-Kantorovich mass transfer problem. The goal is to find the optimal
mapping $M$ which minimizes the Kantorovich-Wasserstein distance. The
optimal mapping can be written as $M=\nabla \psi$, where $\psi$ satisfies
the following Monge-Amp\`{e}re equation
$$
det(D^2 \psi(x)) = \frac{I_1(x)}{I_2(\nabla \psi)},
$$
where $I_1$ and $I_2$ are the given images.
Here $det(D^2 \psi(x))$ denotes the determinant of the Hessian of $\psi$.
In $\mathbb{R}^2$, it is equivalent to (\ref{eqn:MAE}).
In this talk, we will present a multigrid method for solving the
Monge-Amp\`{e}re equation. We will discuss the discretization of the
nonlinear equation and the issues of viscosity solutions and monotone
finite difference and finite element schemes. We will then present a
relaxation scheme which is a very slow convergent method as a standalone
solver but it is very effective for reducing high frequency errors. We will
adopt it as a smoother for multigrid and demonstrate its smoothing
properties. Finally, numerical results will be presented to illustrate
the effectiveness of the method.


\end{document}
